%% ---------------------------------------------------------------- tab:inherited
{\footnotesize
\setlength{\tabcolsep}{4pt}
\renewcommand{\arraystretch}{1.15}
\begin{longtable}{@{}p{0.30\linewidth}p{0.17\linewidth}p{0.48\linewidth}@{}}
\caption{\textbf{What this study developed and what it inherited.}}
\label{tab:inherited}\\
\toprule
\textbf{Component} & \textbf{Origin} & \textbf{Role in this paper} \\
\midrule
\endfirsthead
\multicolumn{3}{@{}l}{\footnotesize\emph{Table \thetable\ continued from previous page.}}\\
\toprule
\textbf{Component} & \textbf{Origin} & \textbf{Role in this paper} \\
\midrule
\endhead
\midrule
\multicolumn{3}{r@{}}{\footnotesize\emph{continued on next page}}\\
\endfoot
\bottomrule
\endlastfoot
Eight interpretability methods (attention analysis, residual-stream
spectral geometry, topological screening, sparse-autoencoder atlas,
circuit tracing, exhaustive mapping and steering, manifold discovery,
donor-aware ageing probe) & Author's earlier studies of scGPT and
Geneformer & Inherited. The source papers of four of them (attention
analysis, topological screening, sparse-autoencoder atlas, circuit
tracing) are given to agents in Study~A as the source method. \\
MaxToki-217M and MaxToki-1B & Third party~\cite{maxtoki2026} & Deployment
target. \\
Perturb-seq screens, Tabula Sapiens, the Human Lung Cell Atlas, TRRUST, DoRothEA, STRING, GO, KEGG,
Reactome & Third party & Data and reference sets. \\
Claude Code and Claude models & Third party (Anthropic) & Agents in the
deployment and subjects in Studies~A--C. \\
\addlinespace
Pipeline specifications (eight markdown files) & This study (written
at the start of the deployment by Claude Opus~4.7) & First part of the framework; the treatment in
Study~A's contract arm, where the executor also received the deployed
specification. \\
Ten-pattern audit checklist & This study, distilled from six earlier
peer-review comment sets & Second part of the framework; the treatment in
the checklist arms of Studies~A--C. \\
Deployment runs on MaxToki-217M and the in-session audit & This study &
Source of the error ledger and of the Study~C snapshot. \\
Independent verification, error ledger, deterministic checker & This
study & Evaluation of the deployment. \\
Studies~A, B and C (tasks, items, keys, protocol) & This study &
Controlled evaluation of the framework. \\
Corrected MaxToki-217M analyses (fixed intervention hooks, corrected
input encoding, retrained sparse autoencoders) & This study & Case study
of MaxToki-217M. \\
\end{longtable}
}

%% ---------------------------------------------------------------- tab:checklist
{\footnotesize
\setlength{\tabcolsep}{4pt}
\renewcommand{\arraystretch}{1.15}
\begin{longtable}{@{}p{0.17\linewidth}p{0.26\linewidth}p{0.28\linewidth}p{0.24\linewidth}@{}}
\caption{\textbf{The ten-pattern audit checklist.}}
\label{tab:checklist}\\
\toprule
\textbf{Pattern} & \textbf{Objection} & \textbf{Signature searched for} &
\textbf{Prescribed repair} \\
\midrule
\endfirsthead
\multicolumn{4}{@{}l}{\footnotesize\emph{Table \thetable\ continued from previous page.}}\\
\toprule
\textbf{Pattern} & \textbf{Objection} & \textbf{Signature searched for} &
\textbf{Prescribed repair} \\
\midrule
\endhead
\midrule
\multicolumn{4}{r@{}}{\footnotesize\emph{continued on next page}}\\
\endfoot
\bottomrule
\endlastfoot
\textbf{P1} Weak null or no chance anchor & ``What would chance alone give?''
& A number without its chance value; a null that only shuffles labels; a
$p$-value without a named null. & Report the null mean and spread next to
every load-bearing number; prefer the strictest null. \\
\textbf{P2} No trivial baseline & ``What does a simpler method give?'' &
A performance figure without a comparison to expression-only or
gene-level baselines. & State every capability as a gain over the
simplest comparator. \\
\textbf{P3} Endpoint does not match the claim & ``The metric measures
something else.'' & A regulatory claim resting on a ranking score; the
same database used to select and to evaluate. & State what the endpoint
literally measures; weaken the claim to match. \\
\textbf{P4} Wrong unit of inference & ``Your unit is donors, not
cells.'' & $p$-values over cells for donor-level claims; random
cross-validation splits when rows share a factor. & Name the unit;
group resampling and splits by the shared factor. \\
\textbf{P5} Selection-driven inflation & ``You selected on one property
and tested another.'' & Hand-picked subsets; ranking and testing on the
same data. & Separate selection from evaluation; control the selected
property. \\
\textbf{P6} No positive control & ``Could the pipeline detect a signal at
all?'' & A negative result with no planted or known-true signal. & Run a
synthetic and a real positive control first. \\
\textbf{P7} Generalisation from one condition & ``One cell line, one
architecture.'' & Conclusions from one dataset, tissue or model without
a scope statement. & Replicate on a held-out condition or state the
scope. \\
\textbf{P8} Unknown stability & ``What is the interval? Another seed?'' &
Point estimates without intervals, threshold sweeps or seed checks. &
Bootstrap at the unit of inference; report threshold and seed
sensitivity. \\
\textbf{P9} Causal overreach & ``Say `is consistent with', not
`confirms'.'' & Causal verbs on correlational evidence. & Weaken the verb
or run the intervention. \\
\textbf{P10} Missing reproducibility scaffolding & ``Provide code,
environment and compute.'' & Cited scripts absent; no pinned environment,
seeds or compute record. & Freeze environment and seeds; publish every
script; report compute. \\
\end{longtable}
\begin{flushleft}
Each row is an
objection that reviewers raised repeatedly on the author's earlier
interpretability papers, rewritten so that an agent can search for it.
AUROC is the area under the receiver-operating-characteristic curve (0.5
is chance); a null model is a randomised version of the data that keeps
some structure and destroys the effect under test.
\end{flushleft}
}

%% ---------------------------------------------------------------- tab:checks
{\footnotesize
\setlength{\tabcolsep}{4pt}
\renewcommand{\arraystretch}{1.15}
\begin{longtable}{@{}p{0.22\linewidth}p{0.36\linewidth}p{0.36\linewidth}@{}}
\caption{\textbf{How the framework's requirements were checked.}}
\label{tab:checks}\\
\toprule
\textbf{Requirement} & \textbf{During the deployment} & \textbf{Deterministic checker} \\
\midrule
\endfirsthead
\multicolumn{3}{@{}l}{\footnotesize\emph{Table \thetable\ continued from previous page.}}\\
\toprule
\textbf{Requirement} & \textbf{During the deployment} & \textbf{Deterministic checker} \\
\midrule
\endhead
\midrule
\multicolumn{3}{r@{}}{\footnotesize\emph{continued on next page}}\\
\endfoot
\bottomrule
\endlastfoot
Specification has the ten sections, in order & Instruction to the agent;
not checked & Detects missing or misordered sections. \\
Pinned code and code references resolve & Agent wrote commit hashes; not
checked & Detects unresolvable pins, malformed references, missing paths
and lines. \\
Parameters used equal the specification, or the deviation is justified &
Agent's choice; partly noted in code comments & Detects deviations from
machine-readable defaults and ranges where the run records the value;
cannot judge whether a deviation is acceptable. \\
Null models actually randomise & Agent judgment (audit pattern P1) &
Checks null size only; cannot tell whether a null randomises. \\
Trivial baselines and positive controls exist & Agent judgment (P2, P6) &
Detects whether the specification has a labelled slot; cannot judge
adequacy. \\
Intervals name their method and resampling unit & Agent judgment (P4, P8)
& Points to interval statements without a method; cannot judge
correctness. \\
Wording matches the evidence & Agent judgment (P3, P9) & Lists causal
words for a reader; cannot judge context. \\
Every reported number comes from the run's outputs & Not checked & Searches the stored outputs for
each number and gives the rate of chance matches; cannot judge
meaning. \\
Pre-registered quality gates are computed and not loosened & Two
pipelines computed their own pass/fail flags & Recomputes stored gate
flags and detects uncomputed gates. \\
Seeds, environment, device and checkpoint are recorded & Agent judgment
(P10) & Detects scripts that use random numbers without a seed, and runs
with no device, package-version, checkpoint or run-manifest record;
cannot tell whether every random call uses the seed. \\
Hypothesis retirement in automated loops & Described in the
specifications (retire after two negatives in a row); not coded in any
pipeline. A scripted batch of six hypotheses in the topological screen
coded its own promote, inconclusive or retire label, with a stop after
three retires in a row that was never reached; a manual review later
changed 4 of the 6 labels & Detects whether a decision log exists. \\
Go/no-go decisions on expensive runs & Author's prompts; not logged &
Not checkable from files. \\
\end{longtable}
\begin{flushleft}
``During
the deployment'' describes what was in place when the eight runs and the
in-session audit were carried out. ``Deterministic checker'' describes
what the checker written afterwards can detect by reading files; what it
cannot detect remains a matter of judgment.
\end{flushleft}
}

%% ---------------------------------------------------------------- tab:pipelines
{\footnotesize
\setlength{\tabcolsep}{4pt}
\renewcommand{\arraystretch}{1.15}
\begin{longtable}{@{}p{0.20\linewidth}p{0.28\linewidth}p{0.36\linewidth}r@{}}
\caption{\textbf{The eight pipelines as run on MaxToki-217M.}}
\label{tab:pipelines}\\
\toprule
\textbf{Pipeline} & \textbf{Question} & \textbf{What was run} & \textbf{Log h} \\
\midrule
\endfirsthead
\multicolumn{4}{@{}l}{\footnotesize\emph{Table \thetable\ continued from previous page.}}\\
\toprule
\textbf{Pipeline} & \textbf{Question} & \textbf{What was run} & \textbf{Log h} \\
\midrule
\endhead
\midrule
\multicolumn{4}{r@{}}{\footnotesize\emph{continued on next page}}\\
\endfoot
\bottomrule
\endlastfoot
Residual-stream spectral geometry & How does each layer's representation
of genes change with depth? & Per-layer singular-value geometry,
effective rank and stability across cell samples; 2,000 TS immune cells;
a later cross-model comparison with Geneformer V2-316M. & 3.6 \\
Attention-derived regulatory networks & Do attention maps recover
regulator--target pairs? & Layer-8 attention against TRRUST and against
knockdown responses; four runs (K562, RPE1, Adamson with MaxToki-217M;
K562 with MaxToki-1B). & 8.8 \\
Topological hypothesis screen & Does higher-order geometry of gene pairs
carry regulatory information? & A subset of the catalogue (21 of 141
hypotheses had code); the strictest null at 3 of 12 layer states; a
cross-model comparison with scGPT. & 2.6 \\
SAE atlas & What do sparse features of the residual stream represent? &
TopK SAEs (4,928 features, $k=32$) at all 12 layer states, trained on 500
K562 cells; annotation; a transcription-factor specificity test;
feature patching. & 10.1 \\
SAE circuit tracing & Which feature-to-feature edges carry causal
influence? & Ablation-based edges from 30 features at each of four layers
(200 K562 cells); a CRISPRi direction-of-change test. & 4.2 \\
Exhaustive mapping and steering & Is the computation built from
combinatorial logic, and can features steer cell state? & Edges from
1,000 layer-5 features (20 cells); four feature triplets (50 cells);
steering at five layers (50 TS immune cells). & 7.7 \\
Manifold discovery & Is there a recoverable developmental ordering of
cells? & A pre-set hematopoietic ordering and 18 candidate orderings,
scored with quality gates on cell-group centroids from TS donors. & 15.2 \\
Donor-aware ageing probe & Is donor age encoded? & First stage only; its
gate failed. & 1.5 \\
\end{longtable}
\begin{flushleft}
``Log
hours'' is the span covered by each pipeline's log files, with
overlapping runs counted once; it is laptop wall-clock time, not GPU
time, and can include idle time. Model computation ran on one laptop;
the runs that recorded a device used the Apple-silicon MPS backend (one
also used the CPU). The 12 layer states are the hidden states the model
exposes: the token embedding (layer~0) and the outputs of its 11
transformer blocks (layers 1--11; layer~11 is taken after the final
normalisation). TS, Tabula Sapiens; SAE, sparse autoencoder.
\end{flushleft}
}

%% ---------------------------------------------------------------- tab:tasks
{\footnotesize
\setlength{\tabcolsep}{4pt}
\renewcommand{\arraystretch}{1.15}
\begin{longtable}{@{}p{0.06\linewidth}p{0.30\linewidth}p{0.22\linewidth}p{0.36\linewidth}@{}}
\caption{\textbf{Study A tasks.}}
\label{tab:tasks}\\
\toprule
& \textbf{Question} & \textbf{Key verdict} & \textbf{Main pitfalls (not stated in the brief)} \\
\midrule
\endfirsthead
\toprule
& \textbf{Question} & \textbf{Key verdict} & \textbf{Main pitfalls} \\
\midrule
\endhead
\bottomrule
\endlastfoot
T1 & Do layer-8 attention scores encode TRRUST regulator--target pairs
(RPE1 cells)? & Not supported (``inconclusive'' accepted under stated
conditions) & Gene-level properties of targets; a degree-preserving null
that actually rewires; few transcription factors as the resampling unit;
model-free pair signals. \\
T2 & Do SAE-circuit predictions get the direction of change after a
CRISPRi knockdown right (K562)? & Not supported (``inconclusive''
accepted under stated conditions) & Unbalanced classes (50\% is not
chance); constant-sign predictors; records clustered by silenced gene; a
model-free target-direction baseline. \\
T3 & Are SAE features that respond to a transcription factor's knockdown
specific to its target genes? & Not supported & Rare and long genes
dominate feature gene lists; specificity requires comparison with other
factors' targets; threshold search; one factor is not the question. \\
T4 & Is MaxToki-217M's gene-embedding geometry aligned with scGPT's, and
how does that compare with Geneformer V2-316M? & Aligned with scGPT, but
clearly less than with Geneformer & scGPT rows must be looked up by
token id; in-sample canonical correlation has a high chance level; genes
are the resampling unit. \\
\end{longtable}
\begin{flushleft}
Each task's data are a fixed snapshot of
the deployment's outputs, except T2. The deployed circuit edges were made
by an intervention hook that deleted a transformer block, so T2's circuit edges come from a re-run
of the deployed circuit tracing with corrected hooks (S2~Appendix,
Amendment~1). The re-run used the same cells, input encoding, sparse
autoencoders and source features as the deployment. Each answer key
refers to its task's data. T2 and T3 use sparse-autoencoder features trained on the deployment's incorrectly
encoded inputs (see \emph{Errors in the deployment's outputs} in
Results); this does not affect the tasks as analysis exercises, because
their keys were computed on the same data.
\end{flushleft}
}
